%! Polynomial Functions and Rates of Change — Unit 2, Lesson 2.2 — Group Activity (Answer Key)
\documentclass[10pt]{article}
\usepackage{precalculus-article}
\usepackage{precalculus-key}   % pulls in -boxes; do NOT also load -boxes
\newcommand{\TallMath}[1]{$\displaystyle #1\rule[-1.4em]{0pt}{3.2em}$}

\begin{document}

\pageheader{Unit 2, Lesson 2.2}{Group Activity --- Answer Key}

\vspace{0.12in}

\begin{scenariobox}[The Drone Survey]{plumlight}
\small
A survey drone flies over a canyon for six minutes. The polynomial function
$a$ gives its altitude in feet $x$ minutes after takeoff, on the restricted
domain $[0, 6]$. Every tier uses this same table and pre-drawn graph --- no
drawing required.

\smallskip
\begin{center}
\renewcommand{\arraystretch}{1.4}
\begin{tabular}{|c|c|c|c|c|c|c|c|}
\hline
$x$ (minutes) & 0 & 1 & 2 & 3 & 4 & 5 & 6 \\
\hline
$a(x)$ (feet) & 200 & 400 & 240 & 100 & 210 & 380 & 500 \\
\hline
\end{tabular}
\end{center}

\begin{center}
\begin{tikzpicture}[x=0.9cm, y=0.62cm]
  \draw[linegray, very thin, step=1] (0,0) grid (6,5);
  \draw[->, thick] (0,0) -- (6.6,0) node[below right] {\small $x$ (min)};
  \draw[->, thick] (0,0) -- (0,5.6) node[above] {\small $a(x)$ (ft)};
  \foreach \x in {1,2,3,4,5,6} \node[below] at (\x,0) {\small \x};
  \node[left] at (0,1) {\small 100};
  \node[left] at (0,2) {\small 200};
  \node[left] at (0,3) {\small 300};
  \node[left] at (0,4) {\small 400};
  \node[left] at (0,5) {\small 500};
  \draw[very thick, plum] plot [smooth] coordinates
    {(0,2) (0.5,3.3) (1,4) (1.5,3.4) (2,2.4) (2.5,1.4) (3,1)
     (3.5,1.3) (4,2.1) (4.5,3.0) (5,3.8) (5.5,4.5) (6,5)};
  \fill[goldacc] (0,2) circle (2.7pt); \node[below right] at (0.05,2) {\small $P$};
  \fill[goldacc] (1,4) circle (2.7pt); \node[above] at (1,4.15) {\small $Q$};
  \fill[goldacc] (3,1) circle (2.7pt); \node[below] at (3,0.9) {\small $R$};
  \fill[goldacc] (6,5) circle (2.7pt); \node[left] at (5.85,4.9) {\small $S$};
\end{tikzpicture}
\end{center}
\end{scenariobox}

\vspace{0.1in}

\boxguard
\begin{tcolorbox}[colback=white, colframe=black!40,
    title=\textbf{Tier R --- Remediate},
    fonttitle=\bfseries, arc=2mm, left=3mm, right=3mm, top=2mm, bottom=2mm]

Read every answer off the labeled graph.

\begin{enumerate}[itemsep=6pt]
  \item The drone is \textbf{climbing} (the graph increases) on
        \ans{$(0,1)$} and on \ans{$(3,6)$} , and \textbf{descending} on
        \ans{$(1,3)$} .

  \item Point $Q$ is a local \ans{maximum} . Its altitude is \ans{400}
        feet, reached at $x = $ \ans{1} minutes.

  \item Point $R$ is a local \ans{minimum} . Its altitude is \ans{100}
        feet, reached at $x = $ \ans{3} minutes.

  \item The domain stops at $x = 0$ and $x = 6$, and both endpoints are
        included. Point $P$ is a local \ans{minimum} , and point $S$ is a
        local \ans{maximum} .

  \item Over the whole flight, the \textbf{global maximum} altitude is
        \ans{500} feet and the \textbf{global minimum} altitude is
        \ans{100} feet.
\end{enumerate}
\end{tcolorbox}

\vspace{0.1in}

\boxguard
\begin{tcolorbox}[colback=white, colframe=black!40,
    title=\textbf{Tier A --- Approaching Proficiency},
    fonttitle=\bfseries, arc=2mm, left=3mm, right=3mm, top=2mm, bottom=2mm]

Work from the \textbf{table} --- cover the graph until question 5.

\begin{enumerate}[itemsep=6pt]
  \item Complete the AROC row (feet per minute).

        \smallskip
        \renewcommand{\arraystretch}{1.4}
        \begin{tabular}{|c|c|c|c|c|c|c|}
        \hline
        \textbf{Interval} & $[0,1]$ & $[1,2]$ & $[2,3]$ & $[3,4]$ & $[4,5]$ & $[5,6]$ \\
        \hline
        \textbf{AROC} & \ans{200} & \ans{$-160$} & \ans{$-140$} & \ans{110} & \ans{170} & \ans{120} \\
        \hline
        \end{tabular}

  \item The AROC changes from \textbf{positive to negative} between the
        intervals $[0,1]$ and $[1,2]$, so a local \ans{maximum} must occur at
        $x = $ \ans{1} .

  \item The AROC changes from \textbf{negative to positive} between the
        intervals $[2,3]$ and $[3,4]$, so a local \ans{minimum} must occur at
        $x = $ \ans{3} .

  \item Now watch the AROCs themselves rather than their signs.

        \begin{enumerate}[label=(\alph*), itemsep=4pt]
          \item The AROCs are \textbf{smallest} over the interval
                \ans{$[1,2]$} , so the rate of change turns around near
                $x = $ \ans{2} --- a point of inflection, where the
                graph switches from concave \ans{down} to concave
                \ans{up} .
          \item The AROCs are \textbf{largest} over the interval
                \ans{$[4,5]$} , so a second point of inflection lies near
                $x = $ \ans{5} .
        \end{enumerate}

  \item Uncover the graph. In one sentence, say why an inflection point is
        \textit{not} the same thing as a local maximum or minimum.

        \vspace{0.05in}
        \ansline{At a local max or min the AROC changes sign --- the drone stops climbing and starts}
        \ansline{descending. At an inflection point the AROC only stops growing and starts shrinking.}
\end{enumerate}
\end{tcolorbox}

\vspace{0.1in}

\boxguard
\begin{tcolorbox}[colback=white, colframe=black!40,
    title=\textbf{Tier E --- Extension},
    fonttitle=\bfseries, arc=2mm, left=3mm, right=3mm, top=2mm, bottom=2mm]

\begin{enumerate}[itemsep=6pt]
  \item A polynomial of degree $n$ can turn around at most $n - 1$ times.
        Fill in the greatest possible number of local extrema.

        \smallskip
        \renewcommand{\arraystretch}{1.4}
        \begin{tabular}{|c|c|c|c|}
        \hline
        \textbf{Degree} & 3 & 4 & 6 \\
        \hline
        \textbf{At most how many local extrema?} & \ans{2} & \ans{3} & \ans{5} \\
        \hline
        \end{tabular}

  \item A polynomial $p$ has exactly four distinct real zeros.

        \begin{enumerate}[label=(\alph*), itemsep=4pt]
          \item $p$ must have at least \ans{3} local extrema. Justify
                using the rule about zeros.

                \vspace{0.05in}
                \ansline{Four zeros make three consecutive pairs, and each pair forces at least one turn.}
          \item Combine that with question 1: the smallest degree $p$ could
                have is \ans{4} .
        \end{enumerate}

  \item Consider the drone's model $a$ over \textit{all} real numbers rather
        than only $[0,6]$.

        \begin{enumerate}[label=(\alph*), itemsep=4pt]
          \item If $a$ has \textbf{even} degree, must it have a global
                maximum or a global minimum? \ans{Yes} \; Explain.

                \vspace{0.05in}
                \ansline{Both ends run the same direction, so the outputs are bounded on one side.}
          \item If $a$ has \textbf{odd} degree, which of the two does it
                have? \ans{Neither one} \; Explain using end behavior.

                \vspace{0.05in}
                \ansline{One end goes to $+\infty$ and the other to $-\infty$, so no output is greatest or least.}
        \end{enumerate}

  \item \textbf{True or false, with justification:} ``A polynomial with
        exactly one local extremum must have degree 2.''

        \vspace{0.05in}
        \ansline{False. Degree 2 is the smallest that works, but $p(x) = x^4$ also has exactly one local}
        \ansline{extremum (a minimum at $x = 0$). ``At most $n-1$'' is a ceiling, not a count.}
\end{enumerate}
\end{tcolorbox}

\end{document}
